% ==============================================================================
% Section 3: Geometric Framework
% "Endogenous Standpoint as a Low-Curvature Attractor"
% ==============================================================================

\section{Geometric Framework}
\label{sec:geometric-framework}

This section develops the abstract geometric framework for endogenous standpoint.
All definitions are stated independently of any specific architectural instantiation;
the transformer-based realization is deferred to Section~\ref{sec:instantiation}.
We work in the category of finite-dimensional real vector spaces with smooth maps.

% ------------------------------------------------------------------------------
\subsection{Event Space}
\label{sec:event-space-abstract}

\begin{definition}[Event Space]
\label{def:event-space}
An \textbf{event space} is a triple $\mathcal{X} = (V, E, T)$ where:
\begin{enumerate}
    \item $V$ is a finite set of \textbf{events} (vertices);
    \item $E \subseteq V \times V$ is a set of directed \textbf{causal edges};
    \item $T \subseteq V \times V \times V$ is a set of \textbf{2-simplices} (consistency
    obligations): each $\sigma = \langle x_i, x_j, x_k \rangle \in T$ represents the
    obligation that the standpoint at $x_k$ should cohere with the standpoint at $x_i$,
    mediated by the intervening event $x_j$.
\end{enumerate}
The 2-simplices $T$ encode the minimal consistency structure: for each $\sigma \in T$,
there exist edges $(x_i, x_j), (x_j, x_k) \in E$, and the pair $(x_i, x_k)$ represents
the expected (direct) transport path.
\end{definition}

% ------------------------------------------------------------------------------
\subsection{Standpoint Fiber and Metric}
\label{sec:fiber-abstract}

\begin{definition}[Standpoint Fiber]
\label{def:fiber}
Let $K$ be a finite index set. The \textbf{standpoint fiber} at event $x \in V$ is
\[
    E_x = \bigoplus_{k \in K} W_k,
\]
where each $W_k$ is a finite-dimensional real vector subspace with $\dim(W_k) = d_k$,
and $\bigoplus$ denotes the ordinary (not necessarily orthogonal) direct sum. We require
$\sum_{k \in K} d_k = d$ for a fixed ambient dimension $d$, so that $E_x \cong
\mathbb{R}^d$.
\end{definition}

\begin{definition}[Fiber Metric]
\label{def:fiber-metric}
The \textbf{canonical fiber metric} on $E_x$ is
\[
    g_F = \sum_{k \in K} \alpha_k\, \langle \cdot, \cdot \rangle_{W_k},
\]
where $\alpha_k > 0$ are \textbf{layer weights} and $\langle \cdot, \cdot \rangle_{W_k}$
is the standard Euclidean inner product restricted to $W_k$. When all $\alpha_k = 1$,
$g_F$ coincides with the standard Euclidean inner product on $\mathbb{R}^d$.

With respect to $g_F$, the decomposition $E_x = \bigoplus_k W_k$ is orthogonal: for
$h_k \in W_k$ and $h_l \in W_l$ with $k \neq l$, $g_F(h_k, h_l) = 0$. We write
$\Pi_k : E_x \to W_k$ for the $g_F$-orthogonal projection onto $W_k$, which satisfies:
\begin{enumerate}
    \item $\Pi_k^2 = \Pi_k$ (idempotent);
    \item $\Pi_k \Pi_l = 0$ for $k \neq l$ (mutually annihilating);
    \item $\sum_{k \in K} \Pi_k = \mathrm{id}_{E_x}$ (resolution of identity).
\end{enumerate}
\end{definition}

% ------------------------------------------------------------------------------
\subsection{Standpoint State and Structure Group}
\label{sec:state-group}

\begin{definition}[Standpoint State]
\label{def:standpoint-state}
The \textbf{standpoint state} at event $x$ is an element $\Psi_x \in E_x$. Its
\textbf{layer components} are $\psi_k(x) = \Pi_k \Psi_x \in W_k$ for each $k \in K$.
We write $\Psi_x = (\psi_k(x))_{k \in K}$ for the tuple of layer components.
\end{definition}

\begin{definition}[Structure Group]
\label{def:structure-group}
The \textbf{structure group} is
\[
    G = \prod_{k \in K} O(d_k),
\]
where $O(d_k)$ is the orthogonal group of $W_k$. An element $g = (g_k)_{k \in K} \in G$
acts on $E_x$ by
\[
    g \cdot h = \sum_{k \in K} g_k\, \Pi_k\, h.
\]
This action preserves the fiber decomposition: each $g \in G$ maps $W_k$ to itself (i.e.,
$g \cdot h \in W_k$ for $h \in W_k$). The fiber metric $g_F$ is $G$-invariant:
$g_F(g \cdot h, g \cdot h') = g_F(h, h')$ for all $h, h' \in E_x$ and $g \in G$.
\end{definition}

% ------------------------------------------------------------------------------
\subsection{Connection and Curvature}
\label{sec:connection-abstract}

\begin{definition}[Connection]
\label{def:connection}
A \textbf{discrete connection} on the standpoint bundle assigns to each edge $(x_i, x_j)
\in E$ a \textbf{transport operator} $U_{ij} : E_{x_i} \to E_{x_j}$ that is linear and
block-preserving: $U_{ij}(W_k) \subseteq W_k$ for all $k \in K$. The transport of a
standpoint state from $x_i$ to $x_j$ is $\Psi_{x_j} = U_{ij} \Psi_{x_i}$.

A \textbf{gauge transformation} replaces $U_{ij}$ by $g_j\, U_{ij}\, g_i^{-1}$ for
$g_i, g_j \in G$.
\end{definition}

\begin{definition}[Curvature]
\label{def:curvature-abstract}
For each 2-simplex $\sigma = \langle x_i, x_j, x_k \rangle \in T$, the \textbf{curvature}
is
\[
    F_{ijk} = U_{ij} \cdot U_{jk} \cdot (U_{ik}^{\exp})^{-1} \in \operatorname{GL}(d,
    \mathbb{R}),
\]
where $U_{ik}^{\exp}$ is the \textbf{expected transport} from $x_i$ to $x_k$ (the
baseline transport estimated from standpoint-preserving conversations). The curvature
measures the path-dependence of transport: whether the composite path through $x_j$ yields
the same result as the expected direct path.

The \textbf{block-specific curvature} for layer $k$ is
\[
    [F_{ijk}]_k = \Pi_k\, F_{ijk}\, \Pi_k, \qquad
    \|F_{ijk}\|_k = \|[F_{ijk}]_k - \Pi_k\|_F.
\]
\end{definition}

% ------------------------------------------------------------------------------
\subsection{Gauge Invariance}
\label{sec:gauge-abstract}

\begin{proposition}[Gauge Invariance of Block-Specific Curvature]
\label{prop:gauge-abstract}
Under the gauge transformation $U_{ij} \mapsto g_j\, U_{ij}\, g_i^{-1}$ with $g_i, g_j
\in G$, the block-specific curvature norms are invariant:
\[
    \big\|[F^g_{ijk}]_k - \Pi_k\big\|_F = \big\|[F_{ijk}]_k - \Pi_k\big\|_F
    \qquad \forall\, k \in K.
\]
\end{proposition}

\begin{proof}
Under the gauge transformation, the curvature transforms by conjugation:
$F_{ijk} \mapsto g_k\, F_{ijk}\, g_k^{-1}$ (since the intermediate $g_j$ cancels).
Taking the $k$-th block:
\[
    [F^g_{ijk}]_k = \Pi_k\, (g_k\, F_{ijk}\, g_k^{-1})\, \Pi_k
    = g_k\, (\Pi_k\, F_{ijk}\, \Pi_k)\, g_k^{-1}
    = g_k\, [F_{ijk}]_k\, g_k^{-1},
\]
where the second equality uses $g_k \Pi_k = \Pi_k g_k$ (since $g_k$ preserves $W_k$).
Therefore:
\[
    [F^g_{ijk}]_k - \Pi_k
    = g_k\, [F_{ijk}]_k\, g_k^{-1} - g_k\, \Pi_k\, g_k^{-1}
    = g_k\, ([F_{ijk}]_k - \Pi_k)\, g_k^{-1}.
\]
Since $g_k \in O(d_k)$ and the Frobenius norm is invariant under orthogonal conjugation
($\|Q A Q^\top\|_F = \|A\|_F$ for $Q \in O(n)$), we conclude $\|[F^g_{ijk}]_k -
\Pi_k\|_F = \|[F_{ijk}]_k - \Pi_k\|_F$.
\end{proof}

% ------------------------------------------------------------------------------
\subsection{Curvature Decomposition}
\label{sec:decomposition}

\begin{lemma}[Block Decomposition of Total Curvature]
\label{lem:decomposition}
When each transport operator $U_{ij}$ is block-diagonal (i.e., $U_{ij} = \bigoplus_k
U_{ij}^{(k)}$ with $U_{ij}^{(k)} : W_k \to W_k$), the curvature $F_{ijk}$ is also
block-diagonal, and the total curvature decomposes exactly:
\[
    \|F_{ijk} - I_d\|_F^2 = \sum_{k \in K} \big\|[F_{ijk}]_k - \Pi_k\big\|_F^2.
\]
\end{lemma}

\begin{proof}
When $U_{ij}$ is block-diagonal, $F_{ijk} = U_{ij} U_{jk} (U_{ik}^{\exp})^{-1}$ is
block-diagonal (product and inverse of block-diagonal matrices are block-diagonal).
Write $M = F_{ijk} - I_d$. Its blocks are:
\[
    [M]_{kl} = \Pi_k\, M\, \Pi_l =
    \begin{cases}
        [F_{ijk}]_k - \Pi_k & \text{if } k = l, \\
        0 & \text{if } k \neq l.
    \end{cases}
\]
Since the decomposition is orthogonal relative to $g_F$, the Frobenius norm satisfies
$\|M\|_F^2 = \sum_{k,l} \|[M]_{kl}\|_F^2 = \sum_k \|[F_{ijk}]_k - \Pi_k\|_F^2$.
\end{proof}

% ------------------------------------------------------------------------------
\subsection{Low-Curvature Attractor}
\label{sec:attractor}

\begin{definition}[Low-Curvature Attractor]
\label{def:attractor}
For $\epsilon > 0$, the \textbf{low-curvature attractor} is
\[
    \mathcal{A}_{\mathrm{self}}
    = \big\{\Psi \in \prod_{x \in V} E_x : \|F_{ijk}\|_k < \epsilon
    \;\;\forall\, \sigma = \langle x_i, x_j, x_k \rangle \in T,\;
    \forall\, k \in K\big\}.
\]
\end{definition}

\begin{proposition}[Properties of $\mathcal{A}_{\mathrm{self}}$]
\label{prop:attractor}
\hfill
\begin{enumerate}
    \item \textbf{Forward closure.} If $\Psi \in \mathcal{A}_{\mathrm{self}}$ and the
    transport operators satisfy $\|U_{ij}^{(k)} - I_{d_k}\|_F < \epsilon / 3$ for all
    edges and layers, then the transported standpoint remains in
    $\mathcal{A}_{\mathrm{self}}$.

    \item \textbf{Perturbation sensitivity.} If a single transport operator $U_{ij}$ has
    a block $U_{ij}^{(k)}$ with $\|U_{ij}^{(k)} - I_{d_k}\|_F > \epsilon$, then the
    standpoint exits $\mathcal{A}_{\mathrm{self}}$ at the next consistency obligation
    involving edge $(x_i, x_j)$.
\end{enumerate}
\end{proposition}

\begin{proof}
(1) By the triangle inequality, $\|F_{ijk}\|_k = \|[U_{ij} U_{jk} (U_{ik}^{\exp})^{-1}]
_k - \Pi_k\|_F$. When all transport blocks are close to identity, the composite transport
is close to the expected transport, so $\|F_{ijk}\|_k$ remains below $\epsilon$.

(2) If $\|U_{ij}^{(k)} - I_{d_k}\|_F > \epsilon$, then the $k$-th block of the composite
transport $U_{ij} U_{jk}$ deviates from the expected transport by more than $\epsilon$,
so $\|F_{ijk}\|_k > \epsilon$ and $\Psi \notin \mathcal{A}_{\mathrm{self}}$.
\end{proof}
